\scanpage{28}
\unsure{ou|où} $\unsure{\mathfrak S|\mathscr S}_3\times\Gamma$ opère sur $\unsure{\Pi|\pi}_n$ via l'opération
évidente de $\deleted{E}\mathfrak S_3$ sur $\mathbb Z^3/\mathbb Z$, et de $\Gamma$
sur $\mu_n$, via le caractère fondamental.
\[
\Gamma\mathrel{\overset{\chi_n}{\text{---}}}(\mathbb Z/n\mathbb Z)^*\tag*{(10)}
\]
Passant à la limite, on trouve une opération
de $\mathfrak S_3\times\Gamma$ sur
\[
\left\{\begin{aligned}
\Pi_\infty=T^3/T&\simeq T\otimes\mathbb Z^3/\mathbb Z\\
&\simeq T\otimes\widehat{\mathbb Z}^3/\widehat{\mathbb Z}\\
\text{où }\quad T&=\varprojlim\mu_n=\overleftarrow{\mu}_\infty(\overline{\mathbb Q})
\end{aligned}\right.\tag*{(11)}
\]
où $\mathfrak S_3$ opère via ses op. sur $\mathbb Z^3$ \unsure{--|ou|$\to$} $\widehat{\mathbb Z}^3$,
et $\Gamma$ par le caractère fondamental
\[
\Gamma\xrightarrow{\chi_\infty=\chi}\widehat{\mathbb Z}^{\times}\tag*{(12)}
\]
Bien entendu,
\[
\left\{\begin{aligned}
\mathfrak S_3\times\Gamma&=\operatorname{Aut}_{\mathbb Q}C_1\\
(\mathfrak S_3\times\Gamma)\underset{1/2}{\mathbin{.}}\mu_n&\overset{\mathrm{dfn}}{=}\unsure{\mathcal E|\Sigma}_n\simeq\operatorname{Aut}_{\mathbb Q}(C_n)\\
(\mathfrak S_3\times\Gamma)\underset{1/2}{\mathbin{.}}\mu_\infty&\overset{\mathrm{dfn}}{=}\mathcal E_\infty\simeq\operatorname{Aut}_{\mathbb Q}(C_\infty)
\end{aligned}\right.\tag*{(13)}
\]

1.5. En fait, l'hom ext.
\diagram{\begin{tikzcd}[ampersand replacement=\&,column sep=large,row sep=small]
\pi_1(\overline C_1^*)\arrow[r]\arrow[d,equal]\&\Pi_\infty\unsure{^{\leftarrow}}\arrow[d,phantom,"{\rotatebox{90}{$\simeq$}}" description]\\
\pi_1(\mathbb P^1_{\overline{\mathbb Q}}\setminus\{0,1,\infty\})\&T^3/T\simeq T\otimes\mathbb Z^3/\mathbb Z
\end{tikzcd}}
définit un isom
% archive sheet number: 27
